% !TITLE 诗经·卫风·氓
% !AUTHOR 无名氏
% !DYNASTY 先秦
% !GENRE 四言诗
% !ORDER 5

\begin{poem}
氓之蚩蚩，抱布贸丝。匪来贸丝，来即我谋。送子涉淇，至于顿丘。匪我愆期，子无良媒。将子无怒，秋以为期。

乘彼垝垣，以望复关。不见复关，泣涕涟涟。既见复关，载笑载言。尔卜尔筮，体无咎言。以尔车来，以我贿迁。

桑之未落，其叶沃若。于嗟鸠兮，无食桑葚！于嗟女兮，无与士耽！士之耽兮，犹可说也。女之耽兮，不可说也。

桑之落矣，其黄而陨。自我徂尔，三岁食贫。淇水汤汤，渐车帷裳。女也不爽，士贰其行。士也罔极，二三其德。

三岁为妇，靡室劳矣；夙兴夜寐，靡有朝矣。言既遂矣，至于暴矣。兄弟不知，咥其笑矣。静言思之，躬自悼矣。

及尔偕老，老使我怨。淇则有岸，隰则有泮。总角之宴，言笑晏晏。信誓旦旦，不思其反。反是不思，亦已焉哉！
\end{poem}

\begin{appreciation}
《氓》是中国文学史上最早、最完整的一首叙事诗，讲的是一个女子从恋爱、结婚到被抛弃的全过程。

诗中的"氓"（méng），是对普通百姓男子的称呼，略带贬义。他起初"蚩蚩"地笑着，假装来换丝，其实是来向女子求婚。女子虽然心动，却坚持"子无良媒"——没有媒人，我不能答应。这反映了先秦时期对婚姻礼仪的重视。但男子生气了，女子心软，约定"秋以为期"。

第二章写女子热恋中的心情。她爬上倒塌的墙头眺望复关——那是氓居住的地方。看不见他的时候，眼泪哗哗地流；看见他的时候，又笑又说。这种情感的剧烈起伏，把恋爱中的女子写得活灵活现。卜筮的结果也很好，于是她带着嫁妆嫁了过去。

第三章是转折点，也是全诗最著名的段落。"桑之未落，其叶沃若"——桑叶还没落的时候，又绿又润，就像女子年轻的容颜。但斑鸠吃多了桑葚会醉倒，女子沉溺于爱情也会迷失自我。"士之耽兮，犹可说也。女之耽兮，不可说也"——男子沉溺于爱情，还可以脱身；女子一旦陷进去，就再也出不来了。这是三千年前的女性觉醒，是血泪换来的教训。

第四章写婚后生活的贫困和男子的变心。"自我徂尔，三岁食贫"——自从嫁到你家，三年都在过苦日子。"女也不爽，士贰其行"——我没有做错什么，是你变了心。"淇则有岸，隰则有泮"——淇水有岸，湿地有边，可是你的背叛却没有尽头。

最后两章写女子的觉醒和决绝。"及尔偕老，老使我怨"——当初你说要和我一起白头到老，现在想到白头我就怨恨。"总角之宴，言笑晏晏"——回忆童年时的快乐，那时的誓言多么真诚。"信誓旦旦，不思其反"——那些诚恳的誓言，没想到你会违背。"反是不思，亦已焉哉"——既然你违背了，那就算了，到此为止吧。

这首诗的伟大之处，在于它写出了一个女性从天真到成熟的全过程。她没有逆来顺受，而是在痛苦中获得了清醒。"亦已焉哉"四个字，是一种自尊，也是一种力量。
\end{appreciation}
